\chapter{奇异簇}\label{app:singular-varieties}
\SourcePage{430}
虽然复流形是镜像对称的重要组成部分，但在许多情形中，奇异簇也会自然出现。例如，\hyperref[chap:mirror-constructions]{第四章}的 Batyrev 构造，处理的是单纯环面簇中可能有奇点的超曲面。本附录回顾本书将遇到的奇点类型。

本书考虑的簇都是 Cohen--Macaulay 簇，即它们的所有局部环都是 Cohen--Macaulay 环。Cohen--Macaulay 簇 $X$ 的关键性质是，它有一个对偶化层，通常记为 $\omega_X$。当 $X$ 光滑时，$\omega_X$ 就是通常的层 $\Omega_X^d$，其中 $d=\dim(X)$；因此，我们常把一般 Cohen--Macaulay 簇的对偶化层写作 $\widehat\Omega_X^d$。Cohen--Macaulay 簇的良好性质在于，它们在对偶理论中表现良好；详细讨论见 \SourceCite{Oda}。$\widehat\Omega_X^d=\omega_X$ 的另一种记法是 $\mathcal O_X(K_X)$，其中 $K_X$ 是典范除子。注意，一般而言，$K_X$ 只是 Weil 除子。

由标准交换代数可知，对 Cohen--Macaulay 簇 $X$，下列条件等价：对偶化层 $\omega_X$ 是线丛；$K_X$ 是 Cartier 除子；所有局部环都是 Gorenstein 环。此时称 $X$ 是 \emph{Gorenstein 簇}。
\index{Cohen--Macaulay簇@Cohen--Macaulay 簇|hyperpage}\index{对偶化层|hyperpage}\index{Gorenstein簇@Gorenstein 簇|hyperpage}

\section{典范奇点与终端奇点}\label{sec:canonical-terminal-singularities}
若正规簇 $X$ 的某一点附近，存在正整数 $r$ 使 $rK_X$ 为 Cartier 除子，并存在一个局部奇点消解 $f:Y\to X$，使得
\begin{equation}\label{eq:canonical-terminal-discrepancies}
rK_Y=f^*(rK_X)+\sum_i a_iE_i,
\end{equation}
其中求和遍及 $f$ 的全部例外除子 $E_i$，且所有 $a_i\geq0$，就称该点是\emph{典范奇点}。若进一步所有 $a_i>0$，就称该奇点是\emph{终端奇点}。若 $X$ 是 Gorenstein 簇，则上式中可取 $r=1$。

这些奇点类型自然地出现在三维簇和更高维簇的典范模型与极小模型研究中。关于奇点本身的更多背景，见 \SourceCite{Reid1,Reid4}。对于 Calabi--Yau 簇，它们的重要性很明显，因为典范类在 Calabi--Yau 的定义中起突出作用。例如，定义~\ref{def:calabi-yau-variety} 中的极小 Calabi--Yau 簇 $V$，具有 Gorenstein、$\mathbb Q$ 析因的终端奇点。因此，若 $V$ 不光滑，则它的任何消解 $f:W\to V$ 都不再是 Calabi--Yau 簇。这是因为 $V$ 的奇点是终端奇点，所以
\[
K_W=f^*K_V+\sum_i a_iE_i
\]
不平凡：$K_V=0$，而 $a_i>0$。

这些奇点的另一个来源是环面簇。\SourceCite{Reid1} 证明，Gorenstein 环面簇至多具有典范奇点。下面将看到，Gorenstein 轨形也有同样的性质。
\index{典范奇点|hyperpage}\index{终端奇点|hyperpage}

\SourcePage{431}
\section{轨形}\label{sec:orbifolds}
先回顾一个簇是轨形的含义。
\begin{definition}\label{def:orbifold-local-chart}
若 $d$ 维簇 $X$ 的每一点 $p\in X$ 都有一个邻域，在解析意义下等价于 $U/G$ 中原点的一个邻域，其中 $G\subset\operatorname{GL}(d,\mathbb C)$ 是除恒等元外不含复反射的有限子群，$U\subset\mathbb C^d$ 是原点的一个 $G$ 不变邻域，则称 $X$ 为\emph{轨形}。
\end{definition}

在这个定义中，\emph{复反射}是 $\operatorname{GL}(d,\mathbb C)$ 的一个有限阶元素，其 $d-1$ 个特征值等于 $1$。定义~\ref{def:orbifold-local-chart} 中的群 $G$ 称为 $\operatorname{GL}(d,\mathbb C)$ 的一个\emph{小子群}，$(U/G,0)$ 称为 $X$ 在 $p$ 处的\emph{局部坐标片}。由 \SourceCite{Prill} 的一个定理，$G$ 在共轭意义下唯一。

“轨形”与“$V$ 流形”这两个名称可以互换，也就是拟光滑的含义。对于环面簇，这些条件等价于其扇是单纯的，此时群 $G$ 是 Abel 群。因此，\IndexLocation{indexloc:1375:431}单纯环面簇至多有有限 Abel 群商奇点。

一个重要直观是：“在 $\mathbb Q$ 上”，轨形的表现如同流形。例如，紧致轨形的奇异上同调满足 Poincaré 对偶性，但须使用 $\mathbb Q$ 系数。同样，对于完备环面簇，$X$ 光滑时 $H^*(X,\mathbb Z)$ 的组合描述，在 $X$ 单纯时仍适用于 $H^*(X,\mathbb Q)$。

接着研究 Gorenstein 轨形。众所周知，每个 Gorenstein 轨形至多具有典范奇点 \SourceCite{Reid1}。不过，还有下面这个广为流传的结果，表明在终端奇点情形下，Gorenstein 轨形的奇异部分相当小。
\begin{proposition}\label{prop:terminal-gorenstein-orbifold-codimension}
若 Gorenstein 轨形至多具有终端奇点，则其奇异轨迹的余维至少为 $4$。
\end{proposition}
\begin{proof}
可归结到 $\mathbb C^d/G$，其中 $G$ 为小子群。给定阶为 $r$ 的非恒等元 $g\in G$，令 $\zeta$ 为一个本原 $r$ 次单位根，将 $g$ 的特征值写成 $\zeta^{a_i}$，其中 $0\leq a_i<r$。由 \SourceCite{Reid1}，Gorenstein 条件意味着 $\sum_{i=1}^d a_i\equiv0\pmod r$；由 \SourceCite{Reid2}，终端条件意味着 $\sum_{i=1}^d a_i>r$。

根据 \SourceCite{Prill}，$\mathbb C^d/G$ 的奇异轨迹是
\[
\{x\in\mathbb C^d:\text{存在 }g\in G\setminus\{1\}\text{ 使 }g(x)=x\}/G.
\]
若奇异轨迹的余维小于 $4$，则能找到 $g\ne1$，使其特征值经适当重排后为 $\zeta^a,\zeta^b,\zeta^c,1,\ldots,1$。Gorenstein 条件与终端条件给出 $r\mid(a+b+c)$ 和 $a+b+c>r$，因而 $a+b+c=2r$。再对 $g^{-1}$ 应用终端判别条件，很容易得到矛盾。细节留给读者。
\end{proof}

轨形同时也是环面簇时，\SourceCite[定理 2.2.9]{Batyrev4} 已证明命题~\ref{prop:terminal-gorenstein-orbifold-codimension}。在三维情形，立即得到如下推论。
\begin{corollary}\label{cor:terminal-gorenstein-threefold-orbifold-smooth}
至多具有终端奇点的三维 Gorenstein 轨形是光滑的。
\end{corollary}

按照\hyperref[chap:introduction]{第一章}中 Calabi--Yau 的定义，$d$ 维轨形 $V$ 是 Calabi--Yau 簇，当且仅当其典范类平凡，且 $H^i(V,\mathcal O)=0$，$i=1,\ldots,d-1$。这是因为典范类平凡意味着 Gorenstein，而对轨形，Gorenstein 又意味着其奇点为典范奇点。此外，上述
\SourcePage{432}
推论表明，三维极小 Calabi--Yau 轨形是光滑的，因为“极小”意味着奇点至多为终端奇点。

接着定义轨形的子轨形。
\begin{definition}\label{def:suborbifold}
给定轨形 $X$，若子簇 $Y\subset X$ 满足：每一点 $p\in Y$ 处，存在 $X$ 的局部坐标片 $(U/G,0)$，使 $Y$ 在 $U$ 中的逆像在 $0$ 处光滑，则称 $Y$ 为\emph{子轨形}。
\end{definition}

容易看出，轨形的子轨形仍是轨形；反之则不成立。$X$ 的一个子簇，即便本身是轨形，也未必是子轨形。这是因为子轨形的奇点与环境空间的奇点有密切联系。
\index{轨形|hyperpage}\index{轨形!子轨形|hyperpage}\index{复反射|hyperpage}\index{小子群|hyperpage}

\section{轨形上的微分形式}\label{sec:orbifold-differential-forms}
轨形上有很好的微分形式理论。给定 $p\in X$，考虑局部坐标片 $(U/G,0)$。将 $U/G$ 上的 $C^\infty$ 形式定义为 $U$ 上的 $G$ 不变 $C^\infty$ 形式。不同坐标片上定义的形式，有自然的粘合方式，由此可定义 de Rham 上同调群 $H^*_{DR}(X,\mathbb R)$；它们与通常的上同调群同构（见 \SourceCite{Satake}）。

轨形上也可进行 Hodge 理论。轨形上的全纯 $p$ 形式用局部坐标片 $(U/G,0)$ 定义，称为 $X$ 上的 \emph{Zariski $p$ 形式}。它们确定一个层，记为 $\widehat\Omega_X^p$。Zariski 层有如下简单刻画。
\begin{proposition}\label{prop:zariski-forms-smooth-locus-pushforward}
若 $X$ 是轨形，$j:X_0\hookrightarrow X$ 为其光滑轨迹的包含映射，则 $\widehat\Omega_X^p=j_*(\Omega_{X_0}^p)$，其中 $\Omega_{X_0}^p$ 为复流形 $X_0$ 上通常的全纯 $p$ 形式层。
\end{proposition}

轨形 $X$ 是 Cohen--Macaulay 簇；可以证明，若 $d=\dim(X)$，则 $\widehat\Omega_X^d$ 为 $X$ 的\IndexLocation{indexloc:0225:432}对偶化层。因此，这一记号与前节一致。层 $\widehat\Omega_X^p$ 还有许多其他良好性质，包括：
\begin{itemize}
\item 存在微分 $d:\widehat\Omega_X^p\to\widehat\Omega_X^{p+1}$，使 $(\widehat\Omega_X^*,d)$ 成为常值层 $\mathbb C_X$ 的一个消解。
\item 存在自然乘积 $\widehat\Omega_X^p\otimes\widehat\Omega_X^q\to\widehat\Omega_X^{p+q}$，使自然映射
\[
\widehat\Omega_X^p\longrightarrow\mathcal Hom_{\mathcal O_X}
(\widehat\Omega_X^{d-p},\widehat\Omega_X^d)
\]
是同构。
\end{itemize}
也可定义 Hodge 群 $H^{p,q}(X)$，此时 Dolbeault 定理
\[
H^q(X,\widehat\Omega_X^p)\simeq H^{p,q}(X)
\]
仍然成立（见 \SourceCite{Baily1}）。对轨形，通往相交上同调的自然映射
\[
H^*(X,\mathbb Q)\longrightarrow IH^*(X,\mathbb Q)
\]
是同构。由于射影簇的相交同调有自然 Hodge 结构（见 \SourceCite{Saito}），可知当 $X$ 为射影轨形时，$H^k(X,\mathbb C)=\bigoplus_{p+q=k}H^{p,q}(X)$ 有纯 Hodge 结构（但须在 $\mathbb Q$ 上考虑），并满足强 Lefschetz 定理。这对\IndexLocation{indexloc:1081:432}单纯多胞形的组合性质有有趣的推论（见 \SourceCite[第 5.6 节]{Fulton3}）。轨形 Hodge 理论的其他参考资料有 \SourceCite{Steenbrink1,Steenbrink2}。

轨形上的 \emph{Kähler 形式}，是一个实的、光滑且闭的 $(1,1)$ 形式，并且在每一点都正定（在坐标片 $(U/G,0)$ 上，这意味着它到 $U$ 的拉回正定）。轨形上 Kähler 形式的更多细节，见 \SourceCite{AGM1,Baily2}。
\index{Zariski形式@Zariski 形式|hyperpage}\index{轨形!Hodge理论@Hodge 理论|hyperpage}\index{轨形!Kähler形式@Kähler 形式|hyperpage}

\SourcePage{433}
\section{轨形的切层}\label{sec:orbifold-tangent-sheaf}
任意代数簇 $X$ 都有切层 $\Theta_X$。当 $X$ 光滑时，还知道 $\Theta_X$ 与 $\widehat\Omega_X^1$ 对偶，且 $H^1(X,\Theta_X)$ 分类 $X$ 的无穷小形变。$X$ 为轨形时，第一个结论仍然成立。
\begin{proposition}\label{prop:orbifold-tangent-dual-zariski-forms}
若 $X$ 为轨形，则 $\Theta_X\simeq\mathcal Hom_{\mathcal O_X}(\widehat\Omega_X^1,\mathcal O_X)$。
\end{proposition}
\begin{proof}
令 $\Omega_X^1$ 为 $X$ 上 Kähler 一形式层。由 $\Omega_X^1$ 的泛性质，$\Theta_X=\mathcal Hom_{\mathcal O_X}(\Omega_X^1,\mathcal O_X)$。若 $j:X_0\to X$ 为 $X$ 的光滑部分的包含映射，则 \SourceCite[第 128 页]{Oda} 的论证表明，可先限制到 $X_0$ 来计算 $\mathcal Hom_{\mathcal O_X}(-,\mathcal O_X)$。命题由此容易推出。
\end{proof}

至于 $H^1(X,\Theta_X)$，回顾一般簇 $X$ 的无穷小形变由 $\operatorname{Ext}^1(\Omega_X^1,\mathcal O_X)$ 分类。对任意轨形，它可能不同于 $H^1(X,\Theta_X)$；但有一种情形，两者一致。
\begin{proposition}\label{prop:terminal-gorenstein-orbifold-deformations}
若 $X$ 是至多具有终端奇点的 Gorenstein 轨形，则 $H^1(X,\Theta_X)$ 分类 $X$ 的无穷小形变。
\end{proposition}
\begin{proof}
\SourceCite{Schlessinger} 的一个结果表明，若轨形的奇异轨迹的余维至少为三，则其层 $\mathcal Ext^1(\Omega_X^1,\mathcal O_X)$ 为零。由命题~\ref{prop:terminal-gorenstein-orbifold-codimension}，这里满足该条件；再由 $\operatorname{Ext}$ 的局部到整体谱序列，立即得到结论。
\end{proof}

对于至多具有终端\IndexLocation{indexloc:1263:433}环面型奇点且在余维三处为 $\mathbb Q$ 析因的簇（参见 \SourceCite[定理 2.2.9]{Batyrev4}），这一命题也适用，而不必假定它们是轨形。

对于任意轨形 $X$，还可证明 $\Theta_X=j_*\Theta_{X_0}$，并且 $\widehat\Omega_X^1$ 是 $\Omega_X^1$ 的双对偶。若 $X$ 为 Gorenstein 轨形，则同构 $\widehat\Omega_X^{d-1}\simeq\mathcal Hom_{\mathcal O_X}(\widehat\Omega_X^1,\widehat\Omega_X^d)$ 与命题~\ref{prop:orbifold-tangent-dual-zariski-forms} 给出
\begin{equation}\label{eq:orbifold-tangent-top-forms}
\widehat\Omega_X^{d-1}\simeq\Theta_X\otimes\widehat\Omega_X^d.
\end{equation}
可如下将它应用于 Calabi--Yau 簇。
\begin{proposition}\label{prop:minimal-cy-orbifold-deformations}
$d$ 维极小 Calabi--Yau 轨形 $V$ 的无穷小形变，由
\[
H^1(V,\Theta_V)\simeq H^1(V,\widehat\Omega_V^{d-1})
\]
分类。
\end{proposition}
\begin{proof}
极小 Calabi--Yau 轨形具有 Gorenstein 终端奇点，所以由命题~\ref{prop:terminal-gorenstein-orbifold-deformations}，$H^1(V,\Theta_V)$ 分类其无穷小形变。由于 $\widehat\Omega_V^d$ 平凡，再由式~\eqref{eq:orbifold-tangent-top-forms} 即得结论。
\end{proof}
\index{轨形!切层|hyperpage}\index{轨形!无穷小形变|hyperpage}

\section{辛轨形}\label{sec:symplectic-orbifolds}
最后，可在 $C^\infty$ 范畴中定义轨形，并有\emph{轨形\IndexLocation{indexloc:0993:433}微分同胚}的概念。\emph{辛轨形}是带有闭的、非退化的二形式的 $C^\infty$ 轨形（二形式用局部坐标片定义）。此外，辛约化过程在此情形下自然适用。细节见 \SourceCite[II-3.6]{Audin}。
\index{辛轨形|hyperpage}
